% Options for packages loaded elsewhere
\PassOptionsToPackage{unicode}{hyperref}
\PassOptionsToPackage{hyphens}{url}
%
\documentclass[
  11pt,
]{article}
\usepackage{lmodern}
\usepackage{amssymb,amsmath}
\usepackage{ifxetex,ifluatex}
\ifnum 0\ifxetex 1\fi\ifluatex 1\fi=0 % if pdftex
  \usepackage[T1]{fontenc}
  \usepackage[utf8]{inputenc}
  \usepackage{textcomp} % provide euro and other symbols
\else % if luatex or xetex
  \usepackage{unicode-math}
  \defaultfontfeatures{Scale=MatchLowercase}
  \defaultfontfeatures[\rmfamily]{Ligatures=TeX,Scale=1}
\fi
% Use upquote if available, for straight quotes in verbatim environments
\IfFileExists{upquote.sty}{\usepackage{upquote}}{}
\IfFileExists{microtype.sty}{% use microtype if available
  \usepackage[]{microtype}
  \UseMicrotypeSet[protrusion]{basicmath} % disable protrusion for tt fonts
}{}
\makeatletter
\@ifundefined{KOMAClassName}{% if non-KOMA class
  \IfFileExists{parskip.sty}{%
    \usepackage{parskip}
  }{% else
    \setlength{\parindent}{0pt}
    \setlength{\parskip}{6pt plus 2pt minus 1pt}}
}{% if KOMA class
  \KOMAoptions{parskip=half}}
\makeatother
\usepackage{xcolor}
\IfFileExists{xurl.sty}{\usepackage{xurl}}{} % add URL line breaks if available
\IfFileExists{bookmark.sty}{\usepackage{bookmark}}{\usepackage{hyperref}}
\hypersetup{
  hidelinks,
  pdfcreator={LaTeX via pandoc}}
\urlstyle{same} % disable monospaced font for URLs
\usepackage[margin=1in]{geometry}
\setlength{\emergencystretch}{3em} % prevent overfull lines
\providecommand{\tightlist}{%
  \setlength{\itemsep}{0pt}\setlength{\parskip}{0pt}}
\setcounter{secnumdepth}{-\maxdimen} % remove section numbering
\usepackage{mathrsfs}
\usepackage{esint}
\setmathfont{Asana Math}
\setlength{\emergencystretch}{3em}

\author{}
\date{23 September 2026}

\begin{document}

\hypertarget{einstein-defect-bridges-without-a-kuxe4hler-assumption}{%
\section{Einstein-defect bridges without a Kähler
assumption}\label{einstein-defect-bridges-without-a-kuxe4hler-assumption}}

23 September 2026. Two conclusions are distinguished: control of the
\textbf{original metric} uses the scalar-band variance estimate and its
spectral factor; control of a \textbf{conformal Yamabe representative}
needs neither. The metrics in both statements need not initially be
Kähler. The all-metric absolute-energy selection theorem is in the
companion \href{line1_energy_selection.md}{energy-selection paper}, not
repeated here.

\hypertarget{common-hypotheses-and-exact-accounting}{%
\subsection{1. Common hypotheses and exact
accounting}\label{common-hypotheses-and-exact-accounting}}

Let \((M,g)\) be a closed connected oriented smooth four-manifold. A
specified Spin\(^c\) class \(a\) is assumed to have ordinary
Seiberg--Witten solutions for every metric conformal to \(g\). The class
may be selected separately for each metric, but is fixed as its
conformal weight varies. Write \[
 p=(a_g^+)^2>0,\quad A=a^2,\quad T=2\chi+3\tau,\quad V=\operatorname{Vol}(g),
\] \[
 S_2=\int s^2,\quad P=\int|\operatorname{Ric}^{\circ}|^2,\quad
 L=\int(s-\bar s)^2,\quad R_2=P+S_2/4,
\] \[
 D=\frac1{4\pi^2}\int\left(\frac{s^2}{24}+\frac{|W^+|^2}{2}\right)-\frac p2.
 \tag{1}
\] Weyl norms are Hilbert--Schmidt operator norms. Classical monopole
inequalities give \(D\ge0\) and \(B_p:=S_2-32\pi^2p\ge0\). Chern--Weil
then gives the exact identities \[
 R_2=8\pi^2(2p-T)+32\pi^2D,\qquad
 P+B_p/4=8\pi^2(p-T)+32\pi^2D.
 \tag{2}
\] These identities and the scalar inequality are classical inputs, not
new existence results. Normalized weighted derivations are given in
\href{round3_monopole_input.md}{normalized monopole derivation}, Section
6.

\hypertarget{the-original-metric-bridge}{%
\subsection{2. The original-metric
bridge}\label{the-original-metric-bridge}}

Assume \(-S\le s\le-s_0<0\), and put \[
 K=S/\sqrt{s_0},\quad C_0=(1+3S/(2s_0))^2,\quad
 C_{\rm var}=\frac23+\frac{4C_0}{3}+\frac{2K^2}{\lambda_1(g)}.
 \tag{3}
\] Here \(\lambda_1(g)>0\) is the first positive scalar Laplace
eigenvalue. The actual-form remainder proved in
\href{line1_scalar_band_remainder.md}{scalar-band remainder}, equation
(11), gives \(L\le144\pi^2C_{\rm var}D\).

Let \(\kappa=-\bar s/4>0\) and
\(Q_\kappa=\int|\operatorname{Ric}+\kappa g|^2=P+L/4\). \textbf{Theorem
1.} Without assuming any Ricci bound, \[
 Q_\kappa\le8\pi^2(p-T)+(32\pi^2+36\pi^2C_{\rm var})D.
 \tag{4}
\] In particular, define the nonnegative supply \[
 \mathcal P=(p-A)+(A-T)_+,\quad
 U=8\pi^2\mathcal P+(32\pi^2+36\pi^2C_{\rm var})D.
 \tag{5}
\] For the relative integral Einstein defect one has \[
 \boxed{Q_\kappa\le U,\qquad
 \mathcal E(g):=\inf_{\lambda\ne0}
 \frac{\int|\operatorname{Ric}-\lambda g|^2}{\lambda^2V}
 \le\frac{4U}{4\kappa^2V+U}\le\frac{16U}{s_0^2V}.}
 \tag{6}
\] \textbf{Proof.} Drop the nonnegative \(B_p/4\) in (2), and add the
scalar-variance bound divided by four. Also \(p-A=\|a_g^-\|_2^2\ge0\),
hence \(p-T\le\mathcal P\). Direct minimization in \(1/\lambda\) gives
\(\mathcal E=4Q_\kappa/(4\kappa^2V+Q_\kappa)\); \(\kappa\ge s_0/4\)
completes the proof.

For a sequence in a fixed scalar band, sufficient conditions are
\(\mathcal P/V\to0\), \(D/V\to0\), and \(D/(\lambda_1V)\to0\). A uniform
spectral gap suffices for the last condition but is not mandatory. It
cannot simply be deleted. At \(D=\mathcal P=0\), the conclusion is
negative Kähler--Einstein; \(D=0\) alone gives only cscK.
Unequal-curvature hyperbolic products have \(D=0\) and \(p-A>0\).

\hypertarget{supplying-an-actual-form-and-its-joint-analytic-budget}{%
\subsection{3. Supplying an actual form and its joint analytic
budget}\label{supplying-an-actual-form-and-its-joint-analytic-budget}}

Constant-rescale to \(\bar s=-4\), and use the scalar-band parameters of
that normalization. Then \(r:=\fint|\operatorname{Ric}+g|^2\le U/V\).
Suppose \(U/V\le r_0\), \(D/V\le d_0\), and \((-a^2)_+/V\le a_0\).
Sections 6--7 of the scalar-band note supply a smooth actual
weighted-monopole form \(\omega\), for every prescribed \(\eta>0\), with
\[
 \|\omega\|_\infty\le\frac{\sqrt3}{2}K,\quad
 B_\omega:=\fint|\Delta_H\omega|^2
       \le C(S,s_0,8+2r_0,a_0,d_0),
 \tag{7}
\] \[
\begin{aligned}
 A_{\omega,1/2}&=\frac12\left(\frac{|\omega|^2}{3}I-\omega\otimes\omega\right),\\
 d_{1/2}&:=\fint\left(|\operatorname{Ric}+g|^2
       +|W^+-A_{\omega,1/2}|^2+(|\omega|^2-2)^2\right)\\
 &\le\frac{4U}{V}+\frac{16\pi^2(10C_0+8)}3\frac D V+\eta.
\end{aligned}
\tag{8}
\] Consequently \[
 (1+B_\omega)d_{1/2}\le C\left(\frac U V+\frac D V+\eta\right).
 \tag{9}
\] The constant in (9) depends only on the displayed budgets. To check
the inputs, \(\fint|\operatorname{Ric}|^2\le8+2r_0\); the bounded-spinor
Hessian estimate controls \(\Delta_H\omega\) without derivatives of the
conformal weight or \(W^-\). This is an explicit invocation of the
inherited Hessian lemma, not a consequence of finite algebra checks. If
the chosen class additionally satisfies \(a^2\ge T\), Chern--Weil gives
\((-a^2)_+/V\le r_0/(8\pi^2)\), so no independent square-cost bound is
needed in this special case.

Unlike a circular composition, Ricci control in (7)--(9) is an
\textbf{output} of Theorem 1. However, \(1/2\) is the realizable
Kähler--Einstein root in \href{round21_analytic.md}{Round21}; (9) is a
calibrated-form and source-budget statement, not a forbidden-root
contradiction or an exact metric repair.

\hypertarget{a-conformal-bridge-without-scalar-band-or-spectral-assumptions}{%
\subsection{4. A conformal bridge without scalar-band or spectral
assumptions}\label{a-conformal-bridge-without-scalar-band-or-spectral-assumptions}}

Return to only the hypotheses of Section 1. Define \[
 \theta(g)=\frac{p-T+4D(g)}p\ge0.
 \tag{10}
\] Nonnegativity follows from (2). \textbf{Theorem 2.} The conformal
class has a negative constant-scalar-curvature Yamabe representative
\(\widehat g\) satisfying \[
 \boxed{\mathcal E(\widehat g)
       \le\frac{4\theta(g)}{1+\theta(g)}.}
 \tag{11}
\] \textbf{Proof.} The Yamabe theorem supplies a minimizer. Its scalar
curvature cannot be nonnegative because the same conformal Spin\(^c\)
structure has \(p>0\) and an actual solution. Thus its Yamabe constant
\(Y<0\). Normalize \(\widehat g=u^2g\) by \(\int u^4=1\). Then \[
 Y=\int(6|du|^2+su^2)\ge-\|s_-\|_2,
 \quad S_2(\widehat g)=Y^2\le S_2(g).
 \tag{12}
\] The quantities \(p\) and \(\int|W^+|^2\) are conformally invariant,
so \(D(\widehat g)\le D(g)\). Since \(\widehat g\) has constant scalar
curvature, \(\mathcal E(\widehat g)=4P(\widehat g)/R_2(\widehat g)\). By
(2) and \(B_p(\widehat g)\ge0\), this is at most \[
 \frac{4[p-T+4D(\widehat g)]}{2p-T+4D(\widehat g)}.
\] The expression increases with \(D\), its derivative having the sign
of \(p>0\). Replacing \(D(\widehat g)\) by \(D(g)\) proves (11).

Theorem 2 changes the metric conformally; it does not remove the
spectral term from Theorem 1 on the original metric. Neither conclusion
upgrades integral approximation to pointwise Ricci pinching or to an
Einstein metric. Relative to Round 17, the direction is reversed: Ricci
and scalar errors are derived here from actual period/topological data
and LeBrun saturation instead of being supplied as hypotheses.

Sources: \href{https://arxiv.org/abs/math/0003068}{LeBrun, Propositions
4.1--4.2}, the classical Yamabe theorem, the scalar-band remainder cited
above, and the \href{round3_monopole_input.md}{normalized monopole
inequalities}. The \href{line1_ricci_remainder_check.py}{algebra
checker} checks finite algebra, including the companion energy-selection
identities; it does not certify these PDE inputs.

\end{document}
